% !PRINTED 24
\weberpage{24}
Wir führen wieder homogene Koordinaten $x_1$, $x_2$, $x_3$, $x_4$ ein
und nehmen die Gleichungen zweier gegebenen Flächen zweiten
Grades in der Form an:
\weq{3}{\begin{aligned} \varphi(x_1,x_2,x_3,x_4) &= \Sigma a_{ik}x_ix_k, \\ \psi(x_1,x_2,x_3,x_4) &= \Sigma b_{ik}x_ix_k. \end{aligned}}

Die beiden Flächen bestimmen ein \gesperrt{Flächenbüschel} zweiter
Ordnung $\xi\varphi=\eta\psi$. Alle Flächen des Büschels schneiden sich
in einer Raumkurve vierter Ordnung, die wir die \gesperrt{Grundkurve}
nennen. Man erhält dieselbe Kurve und dasselbe Büschel, wenn
man die Funktionen $\varphi$ und $\psi$ durch $\varphi'$, $\psi'$ ersetzt, die aus $\varphi$, $\psi$
mittels der linearen Substitution $\begin{pmatrix} m & n \\ p & q \end{pmatrix}$, also durch
\weq{4}{\begin{aligned} \varphi' &= m\varphi+n\psi, \\ \psi' &= p\varphi+q\psi \end{aligned}}
abgeleitet wird, deren Determinante $r=mq-np$ von Null
verschieden ist. Dadurch ergibt sich die Identität
\weq{5}{\xi\varphi+\eta\psi = \xi'\varphi'+\eta'\psi',}
worin
\weq{6}{\begin{aligned} \xi &= m\xi'+p\eta', \\ \eta &= n\xi'+q\eta' \end{aligned}}
eine lineare Transformation der Variablen $\xi$, $\eta$ darstellt. Außer
dieser \gesperrt{binären Substitution} kommt noch eine lineare quaternäre Substitution der Variablen $x_1$, $x_2$, $x_3$, $x_4$ (Koordinatentransformation) in Betracht, deren Determinante wir $=1$ annehmen
können.

Jede quadratische Form besitzt diesen letzteren Transformationen gegenüber eine \gesperrt{Invariante}, nämlich die \gesperrt{Hessesche
Determinante} (Bd.\,I, § 62, 63, 66), und wir erhalten also als
Invariante der Form (5):
\weq{7}{f(\xi,\eta) = \Sigma \pm(\xi a_{11}+\eta b_{11})(\xi a_{22}+\eta b_{22})(\xi a_{33}+\eta b_{33})(\xi a_{44}+\eta b_{44}).}

Dies ist eine binäre biquadratische Form der Variablen $\xi$, $\eta$,
die wir, entwickelt, so darstellen:
\weq{8}{f(\xi,\eta) = a_0\xi^4+a_1\xi^3\eta+a_2\xi^2\eta^2+a_3\xi\eta^3+a_4\eta^4.}

Die Koeffizienten dieser biquadratischen Form $a_0$, $a_1$, $a_2$, $a_3$, $a_4$
ändern sich nicht bei einer Koordinatentransformation. Sie heißen
daher \gesperrt{simultane Invarianten} des Formenpaares $\varphi$, $\psi$ ($a_0$ und
$a_4$ sind die Determinanten von $\varphi$ und von $\psi$).
